\documentclass[11pt]{article}

\usepackage[margin=1in]{geometry}
\usepackage{amsmath,amssymb,amsthm}
\usepackage{listings}
\usepackage{xcolor}
\usepackage{hyperref}
\usepackage{booktabs}

\theoremstyle{plain}
\newtheorem{theorem}{Theorem}
\newtheorem{definition}{Definition}

\lstset{
  basicstyle=\ttfamily\small,
  keywordstyle=\color{blue},
  commentstyle=\color{gray},
  breaklines=true,
  frame=single,
  xleftmargin=1em,
}

\title{A Formalized Lean~4 Proof of the $m=4$ Directed Hamiltonian\\Torus Decomposition via Stochastic Search}

\author{
  Kevin Nebeker\\
  \texttt{kevin@perqed.com}
}

\date{March 2026}

\begin{document}

\maketitle

\begin{abstract}
We present a machine-checked formal proof of the $m=4$ case of the Directed
Hamiltonian Torus Decomposition problem, verified by the Lean~4 kernel via
the \texttt{decide} tactic. Empirical solutions for even $m$ have been found
computationally by several authors, and Morrison formalized Knuth's odd-$m$
construction in Lean~4, but no machine-checked proof of an even-$m$ case has
been published. We use Simulated Annealing with a non-monotonic cooling
schedule to find a valid decomposition of the $4 \times 4 \times 4$ directed
torus (64 vertices, 192 arcs) in 93 seconds on an Apple~M4 processor. We
then verify the result inside Lean's trusted kernel. We describe the energy
function, a subtlety involving in-degree regularity, and the practical
considerations for scaling this approach to larger values of~$m$.
\end{abstract}

\section{Introduction}

Consider a three-dimensional directed graph $G_m$ with vertex set
$V = \{(i, j, k) : 0 \le i, j, k < m\}$ and directed edges from each vertex
$(i,j,k)$ to $(i{+}1, j, k)$, $(i, j{+}1, k)$, and $(i, j, k{+}1)$, all
modulo~$m$. This graph has $m^3$ vertices and $3m^3$ arcs.

\begin{definition}
A \emph{Directed Hamiltonian Torus Decomposition} of $G_m$ is a partition of
the $3m^3$ arcs into exactly 3 directed Hamiltonian cycles, each of
length~$m^3$.
\end{definition}

Knuth posed this problem in the context of Volume~4B of \emph{The Art of
Computer Programming}~\cite{knuth2026}. In February 2026, Anthropic's Claude
Opus~4.6 found a generalized algebraic construction for all odd~$m$, which
Knuth verified and proved correct. During the same exploration, Claude found
isolated empirical solutions for $m = 4$, $m = 6$, and $m = 8$, but could
not generalize them. In a March~6 revision, Knuth noted that Keston
Aquino-Michaels found an elegant decomposition for even~$m$, though the
construction relies on a splice gadget that requires $m \ge 8$. The small
even cases $m = 4$ and $m = 6$ remain outside the scope of that general
proof.

Morrison formalized Knuth's odd-$m$ proof in Lean~4 on March~4, 2026.
To our knowledge, no machine-checked proof of any even-$m$ case has been
published.

We make the following contributions:
\begin{enumerate}
  \item A machine-checked proof of the $m = 4$ decomposition, verified by
  the {\texttt{decide}} tactic inside Lean's trusted kernel (no external
  compiler in the trusted computing base).
  \item A corrected energy function for the stochastic search that accounts
  for in-degree regularity, avoiding a subtle failure mode in functional
  graphs.
  \item A non-monotonic cooling schedule with energy-proportional
  reheating~\cite{kirkpatrick1983,marinari1992}.
\end{enumerate}

The empirical discovery of $m=4$ solutions is not new. Our contribution is
the formal verification and the description of the pipeline.

\section{Method}

\subsection{State Space}

At each vertex, the three outgoing arcs are assigned to three colors
$\{0, 1, 2\}$. There are $3! = 6$ valid assignments per vertex, giving a
search space of $6^{64} \approx 6.3 \times 10^{49}$ for $m=4$.

We represent a state as an array $\sigma \in \{0, \ldots, 5\}^{64}$, where
$\sigma[v]$ selects which permutation of $\{0,1,2\}$ is applied to the
outgoing directions $(+x, +y, +z)$ at vertex~$v$.

The automorphism group of $G_4$ has order $64 \times 6 \times 6 = 2{,}304$,
comprising spatial translations, axis permutations, and color permutations.
We fix $\sigma[0] = 0$ without loss of generality, breaking the $3! = 6$
color permutations and reducing the search space to $6^{63}$. The remaining
384-fold redundancy from translations and axis permutations is not exploited.

\subsection{Energy Function}

Each color subgraph has out-degree~1 at every vertex by construction. However,
the in-degree is not guaranteed to be~1: a vertex may receive zero or multiple
incoming arcs of a given color from its three axial predecessors. A graph with
out-degree~1 but variable in-degree is a functional graph, whose components
consist of directed cycles with trees feeding into them. A naive cycle-counting
energy function can assign $E = 0$ to such a lasso topology, incorrectly
reporting it as a Hamiltonian cycle.

\begin{definition}
Let $C_c(\sigma)$ be the number of components traced by following successors
in the color-$c$ subgraph, and let $D_c(\sigma)$ be the number of vertices
with in-degree $\ne 1$ for color~$c$. The energy is
\[
  E(\sigma) = \sum_{c=0}^{2} \bigl(C_c(\sigma) - 1\bigr)
            + \sum_{c=0}^{2} D_c(\sigma).
\]
\end{definition}

The in-degree penalty ensures that $E = 0$ requires each color subgraph to
be a single cycle visiting all 64 vertices, i.e., a Hamiltonian cycle. The
energy can be computed in $O(m^3)$ time.

\subsection{Non-Monotonic Cooling}

We use the Metropolis-Hastings algorithm with a single-vertex mutation
operator. A random vertex $v \in \{1, \ldots, 63\}$ is selected and its
permutation changed to one of the five alternatives.

The cooling schedule is $T_k = T_0 \cdot \alpha^k$ with $T_0 = 50$ and
$\alpha = 0.999998$. When the best energy has not improved for $W$ iterations,
the temperature is set to $T_r = \max(1,\; E_{\text{best}}^{2/5})$ and the
window is doubled ($W \leftarrow 2W$). Upon improvement, $W$ resets. This is
a standard non-monotonic cooling schedule~\cite{kirkpatrick1983,marinari1992}
with a power-law reheat temperature that scales with the current difficulty
of the optimization.

\subsection{Comparison with Exact Solvers}

For $m = 4$, the decomposition problem is small enough that exact solvers
are practical. An ILP encoding using Miller-Tucker-Zemlin sub-tour
elimination~\cite{miller1960}, or a pure Boolean SAT encoding using
positional reachability constraints, would solve this instance in
milliseconds. Our stochastic approach is not competitive with exact solvers
at this scale.

The SA pipeline evaluates Hamiltonicity in $O(V)$ time with $V$ bytes of
working memory. Whether this memory advantage becomes relevant at larger
scales (e.g., $m = 8$, $V = 512$) is an empirical question that we have
not yet tested.

\section{Results}

\begin{table}[h]
\centering
\caption{Simulated Annealing runs for $G_4$.}
\label{tab:results}
\begin{tabular}{@{}ccccl@{}}
\toprule
Run & Schedule & Reheat & Best $E$ & Outcome \\
\midrule
1 & $\alpha = 0.99999$ & None & 2 & Froze \\
2 & $\alpha = 0.999998$ & Fixed, $T_r\!=\!5$ & 41 & Random walk \\
3 & $\alpha = 0.999998$ & Fixed, $T_r\!=\!2$ & 0 & Solved (old $E$) \\
\textbf{4} & $\alpha = \mathbf{0.999998}$ & $\mathbf{E^{2/5}}$, \textbf{backoff} & $\mathbf{0}$ & \textbf{Solved, 93s} \\
\bottomrule
\end{tabular}
\end{table}

Run~3 used a cycle-counting energy without the in-degree penalty and is
therefore not sound. Run~4 used the corrected energy function, vertex-0
pinning, and the non-monotonic schedule described above. The solution was
found at iteration 1{,}366{,}834 of restart~12, running at approximately
1.5 million iterations per second on an Apple~M4.

\section{Formal Verification}

The discovered witness is a 64-element array of integers in $\{0,\ldots,5\}$.
We encode this array, the successor function, and the verification
predicate as computable Lean~4 definitions using structural recursion on
natural numbers.

The main theorem is:

\begin{lstlisting}[language=ML]
set_option maxRecDepth 131072 in
set_option maxHeartbeats 4000000 in
theorem claude_cycles_m4_valid :
  isValidDecomposition = true := by decide
\end{lstlisting}

The \texttt{decide} tactic reduces the Boolean expression inside Lean's
kernel and confirms that it equals \texttt{true}. Unlike
\texttt{native\_decide}, which compiles the term to C and trusts the
compiler, \texttt{decide} keeps the entire evaluation within the type
checker. The elevated \texttt{maxRecDepth} and \texttt{maxHeartbeats} are
needed because the verification involves $64 \times 64 \times 3 = 12{,}288$
orbit-membership checks.

As a separate check, an independent TypeScript verifier confirms that each
color subgraph forms a single 64-vertex cycle, that all 192 arcs are
covered without overlap, and that every vertex has in-degree~1 and
out-degree~1 per color.

\subsection{Limitations}

Our verification proves that a specific Boolean predicate evaluates to
\texttt{true} on a hardcoded array. It does not include a formal
mathematical specification of the 3D torus topology or Hamiltonian cycles
in Lean's \texttt{Prop} universe. In particular, the correctness of the
\texttt{successor} function (i.e., that it faithfully encodes the edges of
$G_4$) is trusted code, not a verified theorem.

A fully rigorous formalization would define the torus as a directed graph
over $\mathbb{Z}_4^3$, define Hamiltonian decomposition as a
\texttt{Prop}-valued predicate, and prove a reflection theorem establishing
that our Boolean checker implies the mathematical specification. We leave
this to future work.

\section{Conclusion}

We have given a machine-checked proof that a specific Boolean predicate,
encoding the Hamiltonian decomposition property of the $4 \times 4 \times 4$
directed torus, evaluates to \texttt{true} on a discovered witness. The
proof is verified entirely within Lean's trusted kernel via \texttt{decide}.
The correctness of the encoding (that the Boolean predicate faithfully
captures the graph-theoretic property) is argued informally but not
formally verified.

The Aquino-Michaels construction covers even $m \ge 8$. Formalizing that
general proof, or closing the remaining base case $m = 6$, are natural
next steps.

All source code, the TypeScript stochastic engine, the witness payload, and
the Lean~4 proof are available at \url{https://github.com/bneb/perqed}.

\section*{Acknowledgments}

We thank Donald Knuth for posing the problem, Kim Morrison for the Lean~4
formalization of the odd-$m$ case, and Keston Aquino-Michaels for the
general even-$m$ construction.

\begin{thebibliography}{9}
\bibitem{knuth2026}
D.~E. Knuth, ``Three-dimensional Hamiltonian cycles in a torus,''
March 2026. Available: \url{https://cs.stanford.edu/~knuth/papers/claude.pdf}

\bibitem{kirkpatrick1983}
S.~Kirkpatrick, C.~D.~Gelatt, and M.~P.~Vecchi, ``Optimization by simulated
annealing,'' \emph{Science}, vol.~220, no.~4598, pp.~671--680, 1983.

\bibitem{marinari1992}
E.~Marinari and G.~Parisi, ``Simulated tempering: a new Monte Carlo scheme,''
\emph{Europhysics Letters}, vol.~19, no.~6, pp.~451--458, 1992.

\bibitem{miller1960}
C.~E.~Miller, A.~W.~Tucker, and R.~A.~Zemlin, ``Integer programming
formulation of traveling salesman problems,'' \emph{Journal of the ACM},
vol.~7, no.~4, pp.~326--329, 1960.
\end{thebibliography}

\end{document}
